\documentclass[12pt,letterpaper]{article}

% ── Paquetes ──────────────────────────────────────────────────────────────────
\usepackage[utf8]{inputenc}
\usepackage[T1]{fontenc}
\usepackage[spanish]{babel}
\usepackage{amsmath, amssymb, amsthm}
\usepackage{mathtools}
\usepackage{geometry}
\usepackage{enumitem}
\usepackage{hyperref}
\usepackage{xcolor}
\usepackage{titlesec}
\usepackage{mdframed}

\geometry{margin=2.5cm}

% ── Entornos ──────────────────────────────────────────────────────────────────
\newtheoremstyle{ejercicio}
  {10pt}{10pt}{\normalfont}{0pt}{\bfseries}{.}{.5em}{}
\theoremstyle{ejercicio}
\newtheorem{ejercicio}{Ejercicio}[section]

\newtheoremstyle{solucion}
  {6pt}{6pt}{\normalfont}{0pt}{\itshape}{.}{.5em}{}
\theoremstyle{solucion}
\newtheorem*{solucion}{Solución}

\newmdenv[
  linecolor=gray!60,
  backgroundcolor=gray!6,
  roundcorner=4pt,
  innerleftmargin=10pt,
  innerrightmargin=10pt,
  innertopmargin=8pt,
  innerbottommargin=8pt
]{recuadro}

% ── Macros ────────────────────────────────────────────────────────────────────
\newcommand{\E}{\mathbb{E}}
\newcommand{\Var}{\operatorname{Var}}
\newcommand{\Cov}{\operatorname{Cov}}
\newcommand{\R}{\mathbb{R}}
\newcommand{\N}{\mathbb{N}}
\newcommand{\F}{\mathcal{F}}
\newcommand{\Ft}{\mathcal{F}_t}
\newcommand{\BM}{(B_t)_{t\ge 0}}
\newcommand{\Law}{\mathcal{L}}
\newcommand{\iid}{\overset{\text{iid}}{\sim}}
\newcommand{\dto}{\overset{d}{=}}
\newcommand{\as}{\text{ c.s.}}
\newcommand{\sgn}{\operatorname{sgn}}
\renewcommand{\P}{\mathbb{P}}
\newcommand{\norm}[1]{\left\lVert #1 \right\rVert}

% ── Título ────────────────────────────────────────────────────────────────────
\title{%
  \textbf{Tarea-Examen: Movimiento Browniano}\\[4pt]
  \large Soluciones — Secciones II y III (Ejercicios 1–22)\\
  \normalsize Movimiento Browniano, Martingalas y Aplicaciones\\
  Matemáticas Aplicadas y Computación
}
\author{}
\date{}

% ──────────────────────────────────────────────────────────────────────────────
\begin{document}
\maketitle



% ==============================================================================
\section{Propiedades fundamentales, regularidad y escalamiento (Ejercicios 1–12)}
% ==============================================================================

% ─────────────────────────────────────────────────────────────────────────────
\subsection*{Ejercicio 1 — Covarianza y varianza del incremento}
% ─────────────────────────────────────────────────────────────────────────────

\begin{recuadro}
\textbf{Enunciado.} Demuestre que para $0\le s\le t$,
\[
  \E[B_t B_s]=s=\min\{s,t\},\qquad
  \Var(B_t-B_s)=t-s.
\]
Deduzca la matriz de covarianza de $(B_{t_1},\dots,B_{t_n})$.
\end{recuadro}

\begin{solucion}

\textbf{Paso 1 — Propiedad de control.}
El movimiento browniano posee incrementos independientes y centrados:
\[
  B_t - B_s \perp \F_s,\qquad
  B_t - B_s \sim \mathcal{N}(0,\,t-s).
\]

\textbf{Paso 2 — Cálculo de $\E[B_t B_s]$.}
Escribimos $B_t = B_s + (B_t - B_s)$ y usamos la independencia:
\begin{align*}
  \E[B_t B_s]
  &= \E\bigl[(B_s + (B_t-B_s))\,B_s\bigr]\\
  &= \E[B_s^2] + \E[(B_t-B_s)]\,\E[B_s]\\
  &= s + 0\cdot 0 = s.
\end{align*}
Como $0\le s\le t$, tenemos $\min\{s,t\}=s$, luego $\E[B_t B_s]=\min\{s,t\}$.

\textbf{Paso 3 — Varianza del incremento.}
\[
  \Var(B_t-B_s) = \E[(B_t-B_s)^2] - (\E[B_t-B_s])^2 = (t-s) - 0 = t-s.
\]

\textbf{Paso 4 — Matriz de covarianza.}
Sea $0\le t_1<\cdots<t_n$. Para $i\le j$,
\[
  \Cov(B_{t_i}, B_{t_j}) = \E[B_{t_i}B_{t_j}] = \min\{t_i,t_j\} = t_i.
\]
La matriz de covarianza es
\[
  \Gamma = \bigl(\Gamma_{ij}\bigr)_{1\le i,j\le n},
  \qquad
  \Gamma_{ij} = \min\{t_i,\,t_j\}.
\]
Explícitamente,
\[
  \Gamma =
  \begin{pmatrix}
    t_1 & t_1 & \cdots & t_1 \\
    t_1 & t_2 & \cdots & t_2 \\
    \vdots & & \ddots & \vdots \\
    t_1 & t_2 & \cdots & t_n
  \end{pmatrix}.
\]
Esta matriz es semidefinida positiva, lo cual es consistente con que
$(B_{t_1},\dots,B_{t_n})$ sea un vector gaussiano.
\end{solucion}

% ─────────────────────────────────────────────────────────────────────────────
\subsection*{Ejercicio 2 — Incrementos independientes y ley finito-dimensional}
% ─────────────────────────────────────────────────────────────────────────────

\begin{recuadro}
\textbf{Enunciado.} Sea $0<t_1<\cdots<t_n$. Demuestre que el vector de incrementos
$(B_{t_1},\,B_{t_2}-B_{t_1},\,\ldots,\,B_{t_n}-B_{t_{n-1}})$
tiene componentes gaussianas independientes, y reconstruya la ley finito-dimensional de
$(B_{t_1},\ldots,B_{t_n})$.
\end{recuadro}

\begin{solucion}

\textbf{Paso 1 — Gaussianidad conjunta.}
Por definición, el vector de incrementos
\[
  \Delta = \bigl(B_{t_1},\;B_{t_2}-B_{t_1},\;\ldots,\;B_{t_n}-B_{t_{n-1}}\bigr)
\]
es conjuntamente gaussiano: cada componente es gaussiana y el vector se obtiene por una
transformación lineal del proceso browniano, que es gaussiano.

\textbf{Paso 2 — Medias y covarianzas.}
Cada componente tiene media cero: $\E[\Delta_k]=0$ para todo $k$.
Para $k\ne \ell$ y $k<\ell$, el incremento $B_{t_\ell}-B_{t_{\ell-1}}$ es independiente de
$\F_{t_{\ell-1}}\supset\sigma(\Delta_1,\dots,\Delta_k)$, por lo que
\[
  \Cov(\Delta_k,\Delta_\ell) = 0.
\]

\textbf{Paso 3 — Independencia.}
En el caso gaussiano, covarianza nula implica independencia. Luego las componentes de $\Delta$
son independientes. Más precisamente, $\Delta_k=B_{t_k}-B_{t_{k-1}}\sim\mathcal{N}(0,\,t_k-t_{k-1})$, con $t_0:=0$.

\textbf{Paso 4 — Reconstrucción de la ley finito-dimensional.}
El vector original se obtiene por transformación lineal triangular:
\[
  \begin{pmatrix}B_{t_1}\\ B_{t_2}\\ \vdots\\ B_{t_n}\end{pmatrix}
  = L\,\Delta,
  \qquad
  L =
  \begin{pmatrix}
    1 & 0 & 0 &\cdots & 0\\
    1 & 1 & 0 &\cdots & 0\\
    \vdots &   & \ddots & & \vdots\\
    1 & 1 & \cdots & 1 & 1
  \end{pmatrix}.
\]
Como $L$ es determinista, $(B_{t_1},\dots,B_{t_n})$ es gaussiano con media cero y matriz de
covarianza $\Gamma=L\,\mathrm{diag}(t_1-t_0,\,\ldots,\,t_n-t_{n-1})\,L^\top$,
cuya entrada $(i,j)$ es $\min\{t_i,t_j\}$, consistente con el Ejercicio~1.
Su densidad conjunta es
\[
  f(x_1,\dots,x_n)
  = \prod_{k=1}^{n}
    \frac{1}{\sqrt{2\pi(t_k-t_{k-1})}}
    \exp\!\left(-\frac{(x_k-x_{k-1})^2}{2(t_k-t_{k-1})}\right),
  \quad x_0:=0.
\]
\end{solucion}

% ─────────────────────────────────────────────────────────────────────────────
\subsection*{Ejercicio 3 — Propiedad de escalamiento}
% ─────────────────────────────────────────────────────────────────────────────

\begin{recuadro}
\textbf{Enunciado.} Pruebe que para todo $c>0$, el proceso
$\tilde B_t := c^{-1}B_{c^2 t}$
es un movimiento browniano estándar.
\end{recuadro}

\begin{solucion}

\textbf{Paso 1 — Continuidad.} Las trayectorias de $\tilde B$ son continuas (composición de
funciones continuas) y $\tilde B_0 = c^{-1}B_0 = 0$.

\textbf{Paso 2 — Incrementos gaussianos.}
Para $0\le s\le t$,
\[
  \tilde B_t - \tilde B_s = c^{-1}(B_{c^2 t}-B_{c^2 s}).
\]
Como $B_{c^2 t}-B_{c^2 s}\sim\mathcal{N}(0,\,c^2(t-s))$ (por la definición del MB), se tiene
\[
  \tilde B_t - \tilde B_s \sim \mathcal{N}\!\left(0,\,\frac{c^2(t-s)}{c^2}\right)
  = \mathcal{N}(0,\,t-s).
\]

\textbf{Paso 3 — Independencia de incrementos.}
Para $0\le t_1<\cdots<t_n$, los tiempos reescalados $c^2 t_1<\cdots<c^2 t_n$ son distintos,
y los incrementos $B_{c^2 t_k}-B_{c^2 t_{k-1}}$ son independientes por la propiedad del MB.
Por tanto, los incrementos de $\tilde B$ también son independientes.

\textbf{Ejemplo guía ($c=2$).}
\[
  \tilde B_t = \tfrac{1}{2}B_{4t},\quad
  \tilde B_t-\tilde B_s = \tfrac{1}{2}(B_{4t}-B_{4s})\sim\mathcal{N}(0,t-s).\checkmark
\]

\textbf{Conclusión.} $\tilde B$ satisface las tres propiedades de definición del movimiento
browniano estándar: valor inicial nulo, trayectorias continuas, incrementos independientes
estacionarios gaussianos $\mathcal{N}(0,\Delta t)$.
\end{solucion}

% ─────────────────────────────────────────────────────────────────────────────
\subsection*{Ejercicio 4 — Inversión temporal}
% ─────────────────────────────────────────────────────────────────────────────

\begin{recuadro}
\textbf{Enunciado.} Sea $X_t = t\,B_{1/t}$ para $t>0$ y $X_0:=0$. Demuestre que
$(X_t)_{t\ge 0}$ es un movimiento browniano estándar.
\end{recuadro}

\begin{solucion}

\textbf{Paso 1 — Ley finito-dimensional.}
Sea $0<t_1<\cdots<t_n$. El vector
\[
  (X_{t_1},\dots,X_{t_n}) = (t_1 B_{1/t_1},\dots,t_n B_{1/t_n})
\]
es una transformación lineal de $(B_{1/t_1},\dots,B_{1/t_n})$, que es gaussiano
con medias cero. Su covarianza es
\[
  \Cov(X_s,X_t) = st\,\E\!\left[B_{1/s}B_{1/t}\right]
  = st\cdot\min\!\left\{\tfrac{1}{s},\tfrac{1}{t}\right\}
  = st\cdot\frac{1}{\max\{s,t\}} = \min\{s,t\}.
\]
Luego $(X_{t_1},\dots,X_{t_n})$ tiene la misma distribución que $(B_{t_1},\dots,B_{t_n})$.

\textbf{Paso 2 — Continuidad en $t>0$.} Clara, pues $t\mapsto tB_{1/t}$ es continua para $t>0$.

\textbf{Paso 3 — Continuidad en $t=0$: $X_t\to 0$ c.s. cuando $t\to 0^+$.}
Por la ley del logaritmo iterado del MB,
\[
  \limsup_{u\to\infty}\frac{B_u}{\sqrt{2u\log\log u}} = 1\quad\as
\]
Poniendo $u=1/t$,
\[
  |X_t| = t\,|B_{1/t}| \le t\cdot C\sqrt{\frac{2}{t}\log\log\frac{1}{t}}
         = C\sqrt{2t\log\log\frac{1}{t}} \xrightarrow{t\to 0^+} 0\quad\as
\]

\textbf{Conclusión.} $X$ tiene las mismas leyes finito-dimensionales que el MB, trayectorias
continuas y $X_0=0$; por tanto es un movimiento browniano estándar. La interpretación es que
el proceso $B$ y su inversión temporal $X_t=tB_{1/t}$ poseen la misma ley.
\end{solucion}

% ─────────────────────────────────────────────────────────────────────────────
\subsection*{Ejercicio 5 — El MB no es monótono en ningún intervalo}
% ─────────────────────────────────────────────────────────────────────────────

\begin{recuadro}
\textbf{Enunciado.} Pruebe que con probabilidad uno el MB no es monótono en ningún
intervalo no degenerado.
\end{recuadro}

\begin{solucion}

\textbf{Paso 1 — Reducción a rationales.}
Es suficiente probar que, para cada par de racionales $q_1<q_2$, el evento
$A_{q_1,q_2}:=\{B\text{ es no decreciente en }[q_1,q_2]\}$
tiene probabilidad cero; análogamente para no creciente.
La unión numerable $\bigcup_{q_1<q_2\in\mathbb{Q}}A_{q_1,q_2}$ cubre todos los intervalos racionales.

\textbf{Paso 2 — Cálculo de $\P(A_{q_1,q_2})$.}
Fijemos $[a,b]$ con $a<b$. Para una partición diádica $\{a=s_0<s_1<\cdots<s_n=b\}$
con $n$ subintervalos iguales de longitud $\delta=(b-a)/n$,
\[
  A_{a,b} \subseteq \bigcap_{k=0}^{n-1}\{B_{s_{k+1}}-B_{s_k}\ge 0\}.
\]
Los eventos $\{B_{s_{k+1}}-B_{s_k}\ge 0\}$ son independientes, cada uno con probabilidad
$1/2$. Luego
\[
  \P(A_{a,b}) \le \left(\frac{1}{2}\right)^n \xrightarrow{n\to\infty} 0.
\]

\textbf{Paso 3 — Conclusión.}
$\P\!\left(\exists\,q_1<q_2\in\mathbb{Q}:\;B\text{ monótono en }[q_1,q_2]\right)
\le \sum_{q_1<q_2\in\mathbb{Q}}\P(A_{q_1,q_2}) = 0$.

Todo intervalo no degenerado contiene un subintervalo racional, lo que concluye la prueba.
El MB es, por tanto, c.s. de variación no acotada en todo intervalo, consistente con su
variación cuadrática no nula.
\end{solucion}

% ─────────────────────────────────────────────────────────────────────────────
\subsection*{Ejercicio 6 — Variación cuadrática}
% ─────────────────────────────────────────────────────────────────────────────

\begin{recuadro}
\textbf{Enunciado.} Sea $\pi_n=\{k2^{-n}:0\le k\le n2^n\}$.
Estudie $V_n:=\sum_{k=0}^{n2^n-1}(B_{(k+1)2^{-n}}-B_{k2^{-n}})^2$
y pruebe que $V_n\to n$ en $L^2$ y en probabilidad.
\end{recuadro}

\begin{solucion}

\textbf{Paso 1 — Media de $V_n$.}
Denotemos $\Delta_k = B_{(k+1)2^{-n}}-B_{k2^{-n}}\iid\mathcal{N}(0,2^{-n})$.
\[
  \E[V_n] = \sum_{k=0}^{n2^n-1}\E[\Delta_k^2]
           = n2^n\cdot 2^{-n} = n.
\]

\textbf{Paso 2 — Varianza de $V_n$.}
Como los incrementos son independientes y $\E[\Delta_k^4]=3(2^{-n})^2$,
\begin{align*}
  \Var(V_n)
  &= \sum_{k=0}^{n2^n-1}\Var(\Delta_k^2)
   = \sum_{k=0}^{n2^n-1}\bigl(\E[\Delta_k^4]-(\E[\Delta_k^2])^2\bigr)\\
  &= n2^n\bigl(3\cdot 2^{-2n} - 2^{-2n}\bigr)
   = n2^n\cdot 2\cdot 2^{-2n}
   = \frac{2n}{2^n} \xrightarrow{n\to\infty} 0.
\end{align*}

\textbf{Paso 3 — Convergencia.}
\[
  \E\bigl[(V_n-n)^2\bigr] = \Var(V_n) = \frac{2n}{2^n}\to 0,
\]
lo que demuestra $V_n\xrightarrow{L^2} n$.
La convergencia en probabilidad se sigue de la de $L^2$:
\[
  \P\!\left(|V_n-n|>\varepsilon\right)\le\frac{\Var(V_n)}{\varepsilon^2}
  = \frac{2n}{\varepsilon^2\cdot 2^n}\to 0.
\]

\textbf{Interpretación.} Este resultado es la versión discreta de $[B]_t=t$: la variación
cuadrática del MB en $[0,n]$ es exactamente $n$.
\end{solucion}

% ─────────────────────────────────────────────────────────────────────────────
\subsection*{Ejercicio 7 — Las trayectorias no son $\alpha$-Hölder para $\alpha>1/2$}
% ─────────────────────────────────────────────────────────────────────────────

\begin{recuadro}
\textbf{Enunciado.} Pruebe que para todo $\alpha>1/2$, c.s. las trayectorias del MB
no son $\alpha$-Hölder en ningún punto.
\end{recuadro}

\begin{solucion}

\textbf{Paso 1 — Definición y reducción.}
Decimos que $B$ es $\alpha$-Hölder en $t_0$ si existe $C>0$ tal que
$|B_t-B_{t_0}|\le C|t-t_0|^\alpha$ para $t$ cercano a $t_0$.
Es suficiente probar que, para cada diádico $t_0=j2^{-m}$, la propiedad falla c.s.

\textbf{Paso 2 — Incrementos diádicos.}
Sea $\alpha>1/2$ y fijemos $\varepsilon>0$ tal que $\alpha>1/2+\varepsilon$.
Para el incremento diádico en el $k$-ésimo subintervalo de nivel $n$,
\[
  \Delta_{n,k} := B_{(k+1)2^{-n}}-B_{k2^{-n}}\sim\mathcal{N}(0,2^{-n}).
\]
Si $B$ fuera $\alpha$-Hölder en algún punto $t_0\in[k2^{-n},(k+1)2^{-n}]$, entonces
$|\Delta_{n,k}|\le C\cdot 2^{-n\alpha}$.

\textbf{Paso 3 — Borel–Cantelli.}
\begin{align*}
  \P\!\left(|\Delta_{n,k}| > 2^{-n\alpha}\right)
  &= 2\,\P\!\left(\mathcal{N}(0,1)>2^{-n(\alpha-1/2)}\right)\\
  &\ge c\cdot 2^{n(\alpha-1/2)}\cdot \exp\!\left(-\tfrac{1}{2}\cdot 2^{-n(2\alpha-1)}\right)
  \ge c\cdot 2^{-n\beta},
\end{align*}
para algún $\beta<1$ cuando $\alpha$ es solo ligeramente mayor que $1/2$.
Más directamente, se puede mostrar que para $\alpha>1/2$,
\[
  \sum_{n=1}^\infty \P\!\left(\max_{0\le k<2^n}|\Delta_{n,k}| > n\cdot 2^{-n\alpha}\right)
  < \infty,
\]
de modo que por Borel--Cantelli, eventualmente el máximo es menor que $n2^{-n\alpha}$.
Sin embargo, si $B$ fuese $\alpha$-Hölder en $t_0$, el incremento en cada escala $n$ que
contiene a $t_0$ satisfaría $|\Delta_{n,k}|\le C\cdot 2^{-n\alpha}$ para todo $n$ suficientemente
grande. En escala $2^{-n}$ se tiene
$\E[|\Delta_{n,k}|^2]=2^{-n}$, mientras que si $\alpha>1/2$
entonces $2^{-n\alpha}\ll 2^{-n/2}$, lo que contradice el comportamiento típico de la gaussiana.

\textbf{Conclusión.}
Para todo $\alpha>1/2$, la probabilidad de que $B$ sea $\alpha$-Hölder en al menos un punto
es cero. En cambio, las trayectorias son $(1/2-\varepsilon)$-Hölder para todo $\varepsilon>0$
(teorema de Kolmogorov–Čentsov).
\end{solucion}

% ─────────────────────────────────────────────────────────────────────────────
\subsection*{Ejercicio 8 — El conjunto de ceros $Z=\{t\ge 0: B_t=0\}$}
% ─────────────────────────────────────────────────────────────────────────────

\begin{recuadro}
\textbf{Enunciado.} Sea $Z:=\{t\ge 0:B_t=0\}$. Demuestre que $Z$ es cerrado y no
acotado c.s. Discuta por qué esto es compatible con la continuidad de las trayectorias.
\end{recuadro}

\begin{solucion}

\textbf{Paso 1 — $Z$ es cerrado c.s.}
$Z = B^{-1}(\{0\})$. Como $t\mapsto B_t$ es continua c.s. y $\{0\}$ es cerrado en $\R$,
la preimagen $Z$ es cerrada c.s.

\textbf{Paso 2 — $Z$ es no acotado c.s.}
Por la ley del logaritmo iterado:
\[
  \limsup_{t\to\infty}\frac{B_t}{\sqrt{2t\log\log t}} = 1\quad\as,
  \qquad
  \liminf_{t\to\infty}\frac{B_t}{\sqrt{2t\log\log t}} = -1\quad\as
\]
Luego $B$ toma valores positivos y negativos infinitamente frecuentes para $t$ grande.
Por la continuidad de las trayectorias y el teorema del valor intermedio, $B$ debe
anularse en un instante arbitrariamente grande. Formalmente: para cada $T>0$,
\[
  \P\!\left(\exists\,t>T:\;B_t=0\right)=1.
\]
Esto implica que $\sup Z=+\infty$ c.s.

\textbf{Paso 3 — Compatibilidad con la continuidad.}
$Z$ es cerrado y no tiene interior (pues el MB no es monótono, Ejercicio~5).
Es un conjunto de Cantor estocástico: cerrado, sin interior, perfecto y no numerable.
La continuidad de $t\mapsto B_t$ implica que entre dos ceros consecutivos el proceso
excursiona estrictamente por el semiplano positivo o negativo, creando la estructura
de excursiones de Itô. No existe contradicción porque un conjunto cerrado de medida
de Lebesgue cero puede ser no acotado (por ejemplo, $\{0,1,2,\dots\}$), y de hecho
la medida de Lebesgue de $Z$ es cero c.s. (por la fórmula de ocupación y el tiempo local).
\end{solucion}

% ─────────────────────────────────────────────────────────────────────────────
\subsection*{Ejercicio 9 — Vectores gaussianos brownianos}
% ─────────────────────────────────────────────────────────────────────────────

\begin{recuadro}
\textbf{Enunciado.} Sean $0\le s_1\le\cdots\le s_n$. Explique por qué
$(B_{s_1},\dots,B_{s_n})$ tiene distribución normal conjunta.
\end{recuadro}

\begin{solucion}

\textbf{Paso 1 — Vector de incrementos independientes.}
Definamos $\Delta_k = B_{s_k}-B_{s_{k-1}}$ con $s_0:=0$ y $\Delta_1:=B_{s_1}$.
Por definición del MB, las variables $\Delta_1,\dots,\Delta_n$ son independientes con
$\Delta_k\sim\mathcal{N}(0,\,s_k-s_{k-1})$; en particular, son gaussianas.

\textbf{Paso 2 — Transformación lineal triangular.}
El vector original se escribe como
\[
  \begin{pmatrix}B_{s_1}\\ B_{s_2}\\ \vdots\\ B_{s_n}\end{pmatrix}
  = L
  \begin{pmatrix}\Delta_1\\ \Delta_2\\ \vdots\\ \Delta_n\end{pmatrix},
  \qquad
  L_{ij} = \mathbf{1}_{j\le i}.
\]
$L$ es triangular inferior con unos en la diagonal y es determinista.

\textbf{Paso 3 — Gaussianidad conjunta.}
Una transformación lineal de un vector gaussiano es gaussiana. Por tanto
$(B_{s_1},\dots,B_{s_n})$ es gaussiano conjunto con media cero y matriz de covarianza
$\Gamma_{ij}=\Cov(B_{s_i},B_{s_j})=\min\{s_i,s_j\}$.

\textbf{Ejemplo guía ($n=2$).}
\[
  (B_s,B_t) = \begin{pmatrix}1 & 0\\ 1 & 1\end{pmatrix}
              \begin{pmatrix}B_s\\ B_t-B_s\end{pmatrix}.
\]
La covarianza es $\E[B_sB_t]=s=\min\{s,t\}$ para $s\le t$. \checkmark
\end{solucion}

% ─────────────────────────────────────────────────────────────────────────────
\subsection*{Ejercicio 10 — Covarianza, mgf e independencia gaussiana}
% ─────────────────────────────────────────────────────────────────────────────

\begin{recuadro}
\textbf{Enunciado.} Sea $\Gamma$ la matriz de covarianza de $(B_{s_1},\dots,B_{s_n})$.
Calcule $\Gamma$ y la función generadora de momentos. Explique por qué covarianza cero
implica independencia en el caso gaussiano.
\end{recuadro}

\begin{solucion}

\textbf{Paso 1 — Cálculo de $\Gamma$.}
Del Ejercicio~1, $\Gamma_{ij}=\min\{s_i,s_j\}$.

\textbf{Paso 2 — Función generadora de momentos (mgf).}
Para $\theta=(\theta_1,\dots,\theta_n)^\top\in\R^n$, siendo el vector gaussiano
centrado con covarianza $\Gamma$:
\[
  f(\theta_1,\dots,\theta_n)
  = \E\!\left[\exp\!\left(\sum_{i=1}^n\theta_i B_{s_i}\right)\right]
  = \exp\!\left(\frac{1}{2}\,\theta^\top\Gamma\,\theta\right)
  = \exp\!\left(\frac{1}{2}\sum_{i,j}\theta_i\,\theta_j\min\{s_i,s_j\}\right).
\]

\textbf{Paso 3 — Independencia gaussiana.}
Si $\Gamma$ es diagonal, es decir, $\Cov(B_{s_i},B_{s_j})=0$ para $i\ne j$, entonces
\[
  \theta^\top\Gamma\,\theta = \sum_{i=1}^n\theta_i^2\,\Gamma_{ii},
\]
y la mgf factoriza:
\[
  f(\theta_1,\dots,\theta_n) = \prod_{i=1}^n\exp\!\left(\frac{\theta_i^2\,\Gamma_{ii}}{2}\right)
  = \prod_{i=1}^n\E\!\left[e^{\theta_i B_{s_i}}\right].
\]
La factorización de la mgf conjunta implica la independencia de $B_{s_1},\dots,B_{s_n}$.
En general, para un vector gaussiano centrado, covarianza nula equivale a independencia.

\textbf{Ejemplo guía ($n=2$).}
\[
  f(\theta_1,\theta_2)
  = \exp\!\left(\frac{\theta_1^2 s_1 + 2\theta_1\theta_2 s_1 + \theta_2^2 s_2}{2}\right)
  \quad(s_1\le s_2).
\]
Si $s_1=0$, entonces $B_{s_1}=0$ es constante e independiente de todo lo demás.
\end{solucion}

% ─────────────────────────────────────────────────────────────────────────────
\subsection*{Ejercicio 11 — Dos brownianos independientes y encuentros}
% ─────────────────────────────────────────────────────────────────────────────

\begin{recuadro}
\textbf{Enunciado.} Sean $X$, $Y$ dos MB estándar unidimensionales independientes.
Muestre que $Z_t:=X_t-Y_t$ es un MB con varianza $2$. Discuta si
$\{t\ge 0:X_t=Y_t\}$ es c.s. infinito.
\end{recuadro}

\begin{solucion}

\textbf{Paso 1 — Ley de $Z$.}
\[
  Z_t - Z_s = (X_t-X_s)-(Y_t-Y_s),\qquad 0\le s\le t.
\]
Los incrementos $X_t-X_s\sim\mathcal{N}(0,t-s)$ e $Y_t-Y_s\sim\mathcal{N}(0,t-s)$ son independientes, luego
\[
  Z_t-Z_s \sim\mathcal{N}(0,2(t-s)).
\]
Los incrementos de $Z$ en intervalos disjuntos son independientes (por la independencia de $X$ e $Y$).
Las trayectorias de $Z$ son continuas y $Z_0=0$.
Luego $Z$ es un MB con parámetro de varianza $\sigma^2=2$, equivalentemente $Z_t=\sqrt{2}\,\tilde B_t$
donde $\tilde B$ es un MB estándar.

\textbf{Paso 2 — El conjunto de encuentros $\{X_t=Y_t\}=\{Z_t=0\}$.}
El conjunto $\{Z_t=0\}=\{t\ge 0:\tilde B_t=0\}$ es el conjunto de ceros de un MB estándar.
Por el Ejercicio~8, este conjunto es c.s. no acotado. Luego hay infinitos tiempos de encuentro
con probabilidad uno:
\[
  \P\!\left(\#\{t\ge 0:X_t=Y_t\}=\infty\right)=1.
\]
Esto es consistente con la recurrencia del MB unidimensional.
\end{solucion}

% ─────────────────────────────────────────────────────────────────────────────
\subsection*{Ejercicio 12 — Cambio de tiempo por niveles en dimensión dos}
% ─────────────────────────────────────────────────────────────────────────────

\begin{recuadro}
\textbf{Enunciado.} Sea $X_t=(X^1_t,X^2_t)$ un MB bidimensional estándar y
$\sigma_t:=\inf\{s\ge 0:(X^2_s)^2=t\}$, $Y_t:=X^1_{\sigma_t}$.
Determine qué propiedades del MB satisface $Y$. Estudie además para qué $\lambda>0$
el proceso $Z_t:=a^{-\lambda}Y_{at}$ tiene la misma ley que $Y$.
\end{recuadro}

\begin{solucion}

\textbf{Paso 1 — Estudio del cambio de tiempo $t\mapsto\sigma_t$.}
El tiempo $\sigma_t$ es el primer instante en que $|X^2_s|=\sqrt{t}$, es decir,
$\sigma_t=\tau_{\sqrt{t}}$ para el MB $X^2$. Por las propiedades del MB,
$\sigma_t$ es un tiempo de paro finito c.s. (pues $X^2$ es recurrente) y la función
$t\mapsto\sigma_t$ es no decreciente.

\textbf{Paso 2 — Propiedades de $Y_t$.}
Como $X^1$ y $X^2$ son independientes:
\begin{itemize}[leftmargin=*]
  \item \textbf{Valor inicial}: $Y_0 = X^1_{\sigma_0} = X^1_0 = 0$.
  \item \textbf{Continuidad}: $Y$ es continuo c.s. (herencia de $X^1$ y de la continuidad de $\sigma_t$).
  \item \textbf{Incrementos}: $Y_t-Y_s = X^1_{\sigma_t}-X^1_{\sigma_s}$. Dado que
        $X^2$ es independiente de $X^1$, y $\sigma_t$ depende solo de $X^2$,
        los incrementos de $Y$ son condicional (e incondicionalmente) independientes
        en intervalos disjuntos de $t$. Sin embargo, la distribución del incremento
        depende de la distribución de $\sigma_t-\sigma_s$, que \emph{no} es $t-s$
        en general. Por tanto, $Y$ \textbf{no} tiene incrementos estacionarios gaussianos
        de varianza $\Delta t$.
\end{itemize}

\textbf{Paso 3 — Escalamiento.}
Para encontrar $\lambda$ tal que $Z_t=a^{-\lambda}Y_{at}\dto Y_t$, usamos que
$Y_{at}=X^1_{\sigma_{at}}$ y $\sigma_{at}\dto a\,\sigma_t$ (pues la densidad de
$\tau_{\sqrt{t}}$ para el MB es la ley de Lévy con escala $t$, cuya distribución escala
como $a\cdot\sigma_t$ al sustituir $t\mapsto at$). Luego por el escalamiento del MB,
$X^1_{\sigma_{at}}\dto\sqrt{a}\,X^1_{\sigma_t}$, de donde
\[
  Z_t = a^{-\lambda}Y_{at} \dto a^{-\lambda}\cdot\sqrt{a}\cdot Y_t = a^{1/2-\lambda}Y_t.
\]
Para que $Z\dto Y$ se requiere $a^{1/2-\lambda}=1$, es decir $\lambda=\tfrac{1}{2}$.

\textbf{Conclusión.} $Y$ es continuo, centrado en cero, con incrementos independientes pero
no estacionarios (en varianza). El exponente de Hurst natural es $\lambda=1/2$.
\end{solucion}

% ==============================================================================
\newpage
\section{Propiedad de Markov, reflexión y distribuciones explícitas (Ejercicios 13–22)}
% ==============================================================================

% ─────────────────────────────────────────────────────────────────────────────
\subsection*{Ejercicio 13 — Propiedad de Markov simple}
% ─────────────────────────────────────────────────────────────────────────────

\begin{recuadro}
\textbf{Enunciado.} Demuestre la propiedad de Markov simple:
\[
  \E[f(B_t)\mid\F_s] = P_{t-s}f(B_s),\quad 0\le s\le t,
\]
para toda $f$ boreliana acotada, donde $P_t f(x)=\int_{\R}f(y)\,p_t(x,y)\,dy$ y
$p_t(x,y)=(2\pi t)^{-1/2}e^{-(y-x)^2/(2t)}$.
\end{recuadro}

\begin{solucion}

\textbf{Paso 1 — Descomposición del incremento.}
Escribimos $B_t = B_s + (B_t-B_s)$. El incremento $\xi:=B_t-B_s\sim\mathcal{N}(0,t-s)$
es independiente de $\F_s$.

\textbf{Paso 2 — Cálculo de la esperanza condicional.}
Para $f$ boreliana acotada:
\begin{align*}
  \E[f(B_t)\mid\F_s]
  &= \E[f(B_s+\xi)\mid\F_s]\\
  &= \int_{\R}f(B_s+z)\,\frac{1}{\sqrt{2\pi(t-s)}}e^{-z^2/(2(t-s))}\,dz
     \qquad(\text{$\xi$ independiente de }\F_s)\\
  &= \int_{\R}f(y)\,\frac{1}{\sqrt{2\pi(t-s)}}e^{-(y-B_s)^2/(2(t-s))}\,dy\\
  &= P_{t-s}f(B_s).
\end{align*}

\textbf{Formulación equivalente (función armónica).}
Si $h$ es armónica en un dominio $D\subset\R^d$ y $\tau_D$ es el tiempo de salida del
MB de $D$, entonces $h(B_{t\wedge\tau_D})$ es una martingala y
$h(x)=\E^x[h(B_{\tau_D})]$, que es una consecuencia directa de la propiedad de Markov.
\end{solucion}

% ─────────────────────────────────────────────────────────────────────────────
\subsection*{Ejercicio 14 — Principio de reflexión y el máximo}
% ─────────────────────────────────────────────────────────────────────────────

\begin{recuadro}
\textbf{Enunciado.} Sea $\tau_a:=\inf\{t\ge 0:B_t=a\}$, $a>0$. Use el principio de
reflexión para probar que
\[
  \P\!\left(\sup_{0\le u\le t}B_u\ge a\right)=2\P(B_t\ge a).
\]
\end{recuadro}

\begin{solucion}

\textbf{Paso 1 — Descomposición del evento.}
\[
  \left\{\sup_{0\le u\le t}B_u\ge a\right\}
  = \{\tau_a\le t\}.
\]
Descomponemos según el signo de $B_t$ al tiempo $t$:
\begin{align*}
  \P(\tau_a\le t)
  &= \P(\tau_a\le t,\;B_t\ge a)+\P(\tau_a\le t,\;B_t<a).
\end{align*}

\textbf{Paso 2 — Principio de reflexión.}
En el evento $\{\tau_a\le t,\;B_t<a\}$, la trayectoria tocó el nivel $a$ en algún instante
$\tau_a\le t$ y luego bajó. Definamos el proceso reflejado
\[
  \hat{B}_u :=
  \begin{cases}
    B_u & u\le\tau_a,\\
    2a - B_u & u>\tau_a.
  \end{cases}
\]
Por la propiedad de Markov fuerte en $\tau_a$ y la simetría $B\dto -B$ del MB,
$\hat{B}$ es también un MB estándar. Tenemos
\[
  \hat{B}_t > a \iff B_t < a \quad\text{en el evento }\{\tau_a\le t\}.
\]
Luego $\P(\tau_a\le t,\;B_t<a)=\P(\tau_a\le t,\;\hat{B}_t>a)=\P(\tau_a\le t,\;B_t>a)$.

\textbf{Paso 3 — Conclusión.}
Dado que $\{B_t\ge a\}\subseteq\{\tau_a\le t\}$ (si $B_t\ge a$ el proceso ya cruzó $a$):
\begin{align*}
  \P(\tau_a\le t)
  &= \P(\tau_a\le t,\;B_t\ge a)+\P(\tau_a\le t,\;B_t<a)\\
  &= \P(B_t\ge a)+\P(B_t>a)\\
  &= 2\P(B_t\ge a),
\end{align*}
donde hemos usado que $\P(B_t=a)=0$ para la continuidad.
\end{solucion}

% ─────────────────────────────────────────────────────────────────────────────
\subsection*{Ejercicio 15 — Densidad del tiempo de llegada $\tau_a$}
% ─────────────────────────────────────────────────────────────────────────────

\begin{recuadro}
\textbf{Enunciado.} Deduzca del ejercicio anterior la densidad
\[
  \P(\tau_a\in dt)=\frac{a}{\sqrt{2\pi t^3}}\,e^{-a^2/(2t)}\,dt,\quad t>0.
\]
\end{recuadro}

\begin{solucion}

\textbf{Paso 1 — Función de distribución de $\tau_a$.}
Del Ejercicio~14,
\[
  \P(\tau_a\le t) = 2\P(B_t\ge a)
  = 2\int_a^\infty \frac{1}{\sqrt{2\pi t}}\,e^{-x^2/(2t)}\,dx
  = 2\left[1-\Phi\!\left(\frac{a}{\sqrt{t}}\right)\right],
\]
donde $\Phi$ es la función de distribución normal estándar.

\textbf{Paso 2 — Derivación de la densidad.}
Derivamos respecto de $t$. Usando la regla de la cadena con $u=a/\sqrt{t}$:
\begin{align*}
  f_{\tau_a}(t)
  &= \frac{d}{dt}\,2\bigl[1-\Phi(a/\sqrt{t})\bigr]
   = -2\,\phi\!\left(\frac{a}{\sqrt{t}}\right)\cdot\frac{d}{dt}\frac{a}{\sqrt{t}}\\
  &= -2\cdot\frac{1}{\sqrt{2\pi}}\,e^{-a^2/(2t)}\cdot\left(-\frac{a}{2t^{3/2}}\right)\\
  &= \frac{a}{\sqrt{2\pi}\,t^{3/2}}\,e^{-a^2/(2t)}
   = \frac{a}{\sqrt{2\pi t^3}}\,e^{-a^2/(2t)}.
\end{align*}

\textbf{Verificación.} Esta es la densidad de una ley de Lévy con parámetros $\mu=0$,
$c=a^2$. En particular, $\E[\tau_a]=+\infty$ (colas pesadas de orden $t^{-3/2}$).
Se puede verificar que $\int_0^\infty f_{\tau_a}(t)\,dt=1$ por el cambio $u=a^2/(2t)$.
\end{solucion}

% ─────────────────────────────────────────────────────────────────────────────
\subsection*{Ejercicio 16 — Distribución del máximo $M_t=\sup_{0\le s\le t}B_s$}
% ─────────────────────────────────────────────────────────────────────────────

\begin{recuadro}
\textbf{Enunciado.} Pruebe que $M_t\dto|B_t|$. Calcule $\E[M_t]$ y $\Var(M_t)$.
\end{recuadro}

\begin{solucion}

\textbf{Paso 1 — Distribución de $M_t$.}
Del Ejercicio~14, para $a>0$:
\[
  \P(M_t\ge a) = \P\!\left(\sup_{s\le t}B_s\ge a\right) = 2\P(B_t\ge a) = \P(|B_t|\ge a),
\]
donde la última igualdad usa $\P(B_t\ge a)=\P(-B_t\ge a)$ (simetría) y
$\P(|B_t|\ge a)=2\P(B_t\ge a)$ (ambas se cancelan). Luego $M_t\dto|B_t|$.

La densidad de $M_t$ es, para $a>0$:
\[
  f_{M_t}(a) = \frac{d}{da}\P(M_t\le a) = \frac{d}{da}[2\Phi(a/\sqrt{t})-1]
  = \frac{2}{\sqrt{2\pi t}}\,e^{-a^2/(2t)},\quad a>0.
\]

\textbf{Paso 2 — Media de $M_t$.}
$M_t\dto|B_t|\sim|\mathcal{N}(0,t)|$, cuya media es
\[
  \E[M_t] = \E[|B_t|] = \sqrt{\frac{2t}{\pi}}.
\]
Verificación directa: $\int_0^\infty a\cdot\frac{2}{\sqrt{2\pi t}}e^{-a^2/(2t)}da
= \sqrt{\frac{2t}{\pi}}$.

\textbf{Paso 3 — Varianza de $M_t$.}
\[
  \E[M_t^2] = \E[B_t^2] = t.
\]
\[
  \Var(M_t) = \E[M_t^2]-(\E[M_t])^2 = t - \frac{2t}{\pi} = t\left(1-\frac{2}{\pi}\right).
\]
\end{solucion}

% ─────────────────────────────────────────────────────────────────────────────
\subsection*{Ejercicio 17 — Probabilidad de ruina $\P(\tau_a<\tau_{-b})$}
% ─────────────────────────────────────────────────────────────────────────────

\begin{recuadro}
\textbf{Enunciado.} Calcule, para $a,b>0$,
$\P(\tau_a<\tau_{-b})$
e interprételo como versión continua del problema de la ruina del jugador.
\end{recuadro}

\begin{solucion}

\textbf{Paso 1 — Planteamiento martingala.}
Sea $\tau=\tau_a\wedge\tau_{-b}$. Como $B_t$ es una martingala y $\tau$ es un tiempo de paro
acotado ($\tau\le\tau_a\wedge\tau_{-b}$, que es finito c.s. porque el MB es recurrente), el
teorema de parada opcional da:
\[
  \E[B_\tau] = B_0 = 0.
\]
En $\{\tau_a<\tau_{-b}\}$, $B_\tau=a$; en $\{\tau_{-b}<\tau_a\}$, $B_\tau=-b$. Luego:
\[
  0 = a\cdot\P(\tau_a<\tau_{-b}) + (-b)\cdot\P(\tau_{-b}<\tau_a).
\]
Como $\P(\tau_{-b}<\tau_a)=1-\P(\tau_a<\tau_{-b})$:
\[
  0 = a\,p - b\,(1-p),\quad p=\P(\tau_a<\tau_{-b}).
\]

\textbf{Paso 2 — Solución.}
\[
  \boxed{\P(\tau_a<\tau_{-b}) = \frac{b}{a+b}.}
\]

\textbf{Interpretación.}
Este es el análogo continuo del problema de la ruina: un jugador que empieza con capital $0$
(entre la ruina $-b$ y la ganancia $a$) tiene probabilidad $b/(a+b)$ de ganar. A diferencia del
caso de la caminata aleatoria con sesgo, aquí el resultado es independiente de la velocidad
(pues el MB no tiene deriva) y solo depende de las distancias $a$ y $b$ a las barreras.
\end{solucion}

% ─────────────────────────────────────────────────────────────────────────────
\subsection*{Ejercicio 18 — Ecuación de Dirichlet y salida del intervalo $(a,b)$}
% ─────────────────────────────────────────────────────────────────────────────

\begin{recuadro}
\textbf{Enunciado.} Sea $g(x):=P^x(\tau_b<\tau_a)$ con $a<x<b$.
Muestre que $g$ satisface una ecuación de Dirichlet y resuélvala.
\end{recuadro}

\begin{solucion}

\textbf{Paso 1 — Ecuación de Dirichlet.}
Sea $\tau=\tau_a\wedge\tau_b$. El proceso $g(B_{t\wedge\tau})$ es una martingala
(pues $g$ debe ser harmónica en $(a,b)$). Por la propiedad de Markov y la fórmula de Itô:
\[
  dg(B_t) = g'(B_t)\,dB_t + \frac{1}{2}g''(B_t)\,dt.
\]
Para que $g(B_{t\wedge\tau})$ sea martingala, el término de deriva debe anularse:
\[
  \frac{1}{2}g''(x) = 0,\quad x\in(a,b),
\]
con condiciones de frontera $g(a)=0$ y $g(b)=1$.

\textbf{Paso 2 — Solución.}
$g''=0$ implica $g(x)=\alpha x+\beta$ (función afín). Aplicando las condiciones de frontera:
\[
  \alpha a+\beta = 0,\quad \alpha b+\beta=1.
\]
Restando: $\alpha(b-a)=1\implies\alpha=\frac{1}{b-a}$. Luego $\beta=-\frac{a}{b-a}$.
\[
  \boxed{g(x) = P^x(\tau_b<\tau_a) = \frac{x-a}{b-a},\quad a\le x\le b.}
\]
Consistentemente, $g(a)=0$, $g(b)=1$ y la fórmula es creciente en $x$.
Como caso especial, con $a=0$ y $x=0$: $g(0)=\P^0(\tau_b<\tau_0)=0$ (reflexión en 0).
\end{solucion}

% ─────────────────────────────────────────────────────────────────────────────
\subsection*{Ejercicio 19 — Distribución del supremo hasta el tiempo de salida de $(-1,1)$}
% ─────────────────────────────────────────────────────────────────────────────

\begin{recuadro}
\textbf{Enunciado.} Sea $T:=\inf\{t\ge 0:|B_t|=1\}$. Calcule la distribución de
$\sup_{0\le s\le T}B_s$.
\end{recuadro}

\begin{solucion}

\textbf{Paso 1 — Identificación del evento.}
Sea $M:=\sup_{0\le s\le T}B_s$. Para $0\le c\le 1$, el evento $\{M\ge c\}$ ocurre si y solo si
el MB llega al nivel $c$ antes de salir de $(-1,1)$, lo cual incluye la posibilidad de llegar
a $1$ (que es la salida superior).

\textbf{Paso 2 — Probabilidad de cruce.}
Para $0\le c\le 1$, el MB parte de $0$ dentro del intervalo $(-1,1)$.
La probabilidad de alcanzar el nivel $c\in(0,1]$ antes de $-1$ es, por el Ejercicio~18
con $a=-1$ y $b=c$:
\[
  P^0(\tau_c<\tau_{-1}) = \frac{0-(-1)}{c-(-1)} = \frac{1}{1+c}.
\]
Dado que $M\ge c$ si y solo si el proceso alcanza $c$ antes de ser absorbido en $-1$
(notar que alcanzar $1$ también implica $M\ge c$ para $c\le 1$), hay que distinguir casos:

\begin{itemize}
  \item Si $c=1$: $\P(M\ge 1)=\P(B_T=1)$.
  \item Para $0<c<1$: $\{M\ge c\}$ ocurre si $\tau_c<\tau_{-1}$, lo cual incluye
        llegar luego a $1$.
\end{itemize}
Más precisamente, $\P(M\ge c)$ es la probabilidad de cruzar $c$ antes que $-1$:
\[
  \P(M\ge c) = P^0(\tau_c<\tau_{-1}) = \frac{1}{1+c},\quad 0<c\le 1.
\]
Verificación: $\P(M\ge 0)=1$ (trivialmente, pues $B_0=0$) y $\P(M\ge 1)=1/2$ (por simetría,
la probabilidad de salir por $1$ es $1/2$).

\textbf{Paso 3 — Densidad.}
\[
  \P(M\le c) = 1 - \frac{1}{1+c},\quad c\in(0,1].
\]
\[
  f_M(c) = \frac{d}{dc}\left(1-\frac{1}{1+c}\right) = \frac{1}{(1+c)^2},\quad c\in(0,1].
\]
\end{solucion}

% ─────────────────────────────────────────────────────────────────────────────
\subsection*{Ejercicio 20 — Ley de $(M_t-B_t, M_t)$ y reflexión}
% ─────────────────────────────────────────────────────────────────────────────

\begin{recuadro}
\textbf{Enunciado.} Demuestre (o establezca una versión parcial) de que la ley de
$(M_t-B_t,\,M_t)$ coincide con la de $(|B_t|,\,L^0_t)$.
\end{recuadro}

\begin{solucion}

\textbf{Versión parcial: identificación marginal de $M_t-B_t$.}

\textbf{Paso 1 — Distribución de $M_t-B_t$.}
Para $x>0$:
\begin{align*}
  \P(M_t-B_t>x)
  &= \P(M_t>B_t+x).
\end{align*}
Condicionando sobre el valor de $B_t=y$:
\[
  \P(M_t-B_t>x) = \int_{-\infty}^\infty \P(M_t>y+x\mid B_t=y)\,f_{B_t}(y)\,dy.
\]
Usando reflexión: $\P(M_t>y+x\mid B_t=y)=\P(M_t>y+x,B_t=y)/\P(B_t=y)$, que por el
principio de reflexión resulta en
\[
  \P(M_t-B_t>x)=\P(|B_t|>x)=2\P(B_t>x),\quad x>0.
\]
Esto muestra $M_t-B_t\dto|B_t|$.

\textbf{Paso 2 — Identidad de Lévy (resultado completo).}
La identidad completa $(M_t-B_t,M_t)\dto(|B_t|,L^0_t)$ es la \textbf{identidad de Lévy},
donde $L^0_t$ es el tiempo local en $0$ hasta $t$. Esta identidad se prueba usando la
fórmula de Tanaka $|B_t|=\int_0^t\text{sgn}(B_s)dB_s+L^0_t$ y mostrando que el par
$(|B_t|,L^0_t)$ satisface la misma ecuación que $(M_t-B_t,M_t)$ vía la reflexión de Skorokhod.
En particular, $L^0_t\dto M_t$ en distribución marginal.
\end{solucion}

% ─────────────────────────────────────────────────────────────────────────────
\subsection*{Ejercicio 21 — Probabilidades de signo: $\P(B_2>0\mid B_1>0)$}
% ─────────────────────────────────────────────────────────────────────────────

\begin{recuadro}
\textbf{Enunciado.} Calcule $\P(B_2>0\mid B_1>0)$ y discuta si $\{B_1>0\}$ y
$\{B_2>0\}$ son independientes.
\end{recuadro}

\begin{solucion}

\textbf{Paso 1 — Descomposición.}
$B_2 = B_1 + (B_2-B_1)$, donde $B_2-B_1\sim\mathcal{N}(0,1)$ es independiente de $B_1$.

\textbf{Paso 2 — Cálculo.}
\begin{align*}
  \P(B_2>0,\,B_1>0)
  &= \P(B_1+(B_2-B_1)>0,\,B_1>0)\\
  &= \int_0^\infty\P(B_2-B_1>-y)\,\frac{1}{\sqrt{2\pi}}e^{-y^2/2}\,dy\\
  &= \int_0^\infty\Phi(y)\,\frac{1}{\sqrt{2\pi}}e^{-y^2/2}\,dy.
\end{align*}
Por simetría y usando que $(B_1,B_2-B_1)$ son independientes estándar, calculamos la
probabilidad de que ambas coordenadas del vector gaussiano $(B_1,B_2)$ sean positivas.
La correlación entre $B_1$ y $B_2$ es $\rho=1/\sqrt{2}$:
\[
  \P(B_1>0,\,B_2>0) = \frac{1}{4}+\frac{1}{2\pi}\arcsin(\rho)
  = \frac{1}{4}+\frac{1}{2\pi}\cdot\frac{\pi}{4}=\frac{1}{4}+\frac{1}{8}=\frac{3}{8}.
\]
Luego,
\[
  \P(B_2>0\mid B_1>0) = \frac{\P(B_2>0,\,B_1>0)}{\P(B_1>0)} = \frac{3/8}{1/2} = \frac{3}{4}.
\]

\textbf{Paso 3 — Independencia.}
$\P(B_1>0)\cdot\P(B_2>0) = \frac{1}{2}\cdot\frac{1}{2}=\frac{1}{4}\ne\frac{3}{8}$.
Por tanto, $\{B_1>0\}$ y $\{B_2>0\}$ \textbf{no son independientes}; están positivamente
correlacionados, pues el MB tiene memoria (su variación cuadrática es finita).
\end{solucion}

% ─────────────────────────────────────────────────────────────────────────────
\subsection*{Ejercicio 22 — Máximo browniano en tiempo 1}
% ─────────────────────────────────────────────────────────────────────────────

\begin{recuadro}
\textbf{Enunciado.} Sea $M:=\max\{B_t:0\le t\le 1\}$. Encuentre la densidad de $M$
y calcule $\E[M]$ y $\Var(M)$.
\end{recuadro}

\begin{solucion}

\textbf{Paso 1 — Función de distribución.}
Para $x>0$, por el principio de reflexión (Ejercicio~14 con $t=1$):
\[
  \P(M\ge x) = 2\P(B_1\ge x).
\]
Luego $\P(M\le x) = 1-2\P(B_1\ge x) = 2\Phi(x)-1$, donde $\Phi$ es la CDF normal estándar.

\textbf{Paso 2 — Densidad.}
\[
  f_M(x) = \frac{d}{dx}[2\Phi(x)-1] = 2\phi(x) = \sqrt{\frac{2}{\pi}}\,e^{-x^2/2},
  \quad x>0.
\]
Esta es la densidad de $|B_1|$, consistente con $M\dto|B_1|$ del Ejercicio~16.

\textbf{Paso 3 — Media.}
\[
  \E[M] = \int_0^\infty x\cdot\sqrt{\frac{2}{\pi}}\,e^{-x^2/2}\,dx
  = \sqrt{\frac{2}{\pi}}\left[-e^{-x^2/2}\right]_0^\infty
  = \sqrt{\frac{2}{\pi}}.
\]

\textbf{Paso 4 — Varianza.}
\[
  \E[M^2] = \int_0^\infty x^2\cdot\sqrt{\frac{2}{\pi}}\,e^{-x^2/2}\,dx
  = \E[B_1^2] = 1.
\]
\[
  \Var(M) = \E[M^2]-(\E[M])^2 = 1-\frac{2}{\pi} = \frac{\pi-2}{\pi}\approx 0.363.
\]
\end{solucion}

% ──────────────────────────────────────────────────────────────────────────────
\newpage
\appendix
\section*{Tabla de resultados principales}
\addcontentsline{toc}{section}{Tabla de resultados principales}

\begin{center}
\renewcommand{\arraystretch}{1.6}
\begin{tabular}{|c|l|}
\hline
\textbf{Objeto} & \textbf{Resultado} \\
\hline
$\Cov(B_s,B_t)$ & $\min\{s,t\}$ \\
$\Var(B_t-B_s)$ & $t-s$ \\
Escala $\tilde B_t=c^{-1}B_{c^2t}$ & MB estándar \\
Inversión $X_t=tB_{1/t}$ & MB estándar \\
$V_n\to$ & $n$ en $L^2$ \\
$[B]_t$ & $t$ (c.s.) \\
$\P(\sup_{s\le t}B_s\ge a)$ & $2\P(B_t\ge a)=2(1-\Phi(a/\sqrt{t}))$ \\
$f_{\tau_a}(t)$ & $\dfrac{a}{\sqrt{2\pi t^3}}e^{-a^2/(2t)}$ \\
$M_t=\sup_{s\le t}B_s$ & $M_t\dto|B_t|$, $\E[M_t]=\sqrt{2t/\pi}$, $\Var=t(1-2/\pi)$ \\
$\P(\tau_a<\tau_{-b})$ & $b/(a+b)$ \\
$P^x(\tau_b<\tau_a)$ & $(x-a)/(b-a)$ \\
$\P(B_2>0|B_1>0)$ & $3/4$ \\
$\E[M]$, $M=\max_{[0,1]}B$ & $\sqrt{2/\pi}$ \\
$\Var(M)$ & $(\pi-2)/\pi$ \\
\hline
\end{tabular}
\end{center}

\bigskip
\noindent\textbf{Convenciones.} $\Phi(x)=\P(\mathcal{N}(0,1)\le x)$,
$\phi(x)=(2\pi)^{-1/2}e^{-x^2/2}$.
Todos los resultados son para $(B_t)_{t\ge 0}$ movimiento browniano estándar partiendo de $0$.

\end{document}
